% Libro 4 · Mundo 2 · Aplicación y analogía (app: pu-aplicacion)
\mundo{Aplicación y analogía}{pu-aplicacion}{Entender una idea de verdad es poder usarla fuera del texto: reconocerla en
un caso nuevo o ver que otra situación funciona de la misma manera. El examen lo pregunta así: «¿cuál de los siguientes
casos ilustra mejor la idea del autor?» o «¿cuál situación es análoga a la descrita?».}{%
\item \textbf{2.1} Aplicar a un caso nuevo
\item \textbf{2.2} Analogías}

\leccion{pu-ap-1}{Aplicar a un caso nuevo}{Juego de llevar una regla del texto a otra situación.}

\begin{reto}
Un texto sostiene que «las personas valoran más lo que les costó esfuerzo conseguir». ¿Cuál caso ilustra mejor esa idea?
\opcionesl{Una persona aprecia más la mesa que armó ella misma que una idéntica comprada armada.}{Una persona compra lo más barato.}{Un niño pierde un juguete que le regalaron.}{Alguien prefiere el transporte rápido.}
\end{reto}

\tecnica{La técnica: extraer el principio}
Para aplicar una idea a un caso nuevo hay que separarla de su ejemplo. Primero se escribe el \clave{principio} en forma
general (sin los nombres ni los detalles del texto); después se busca la opción que tenga \emph{la misma estructura},
aunque hable de otra cosa. Los distractores suelen compartir el \emph{tema} del texto pero no su principio.

\figura{principio}{Del ejemplo se extrae el principio, y el principio se lleva al caso nuevo.}

\begin{regla}
\begin{enumerate}
  \item Escribe la idea como regla general: «cuando X, entonces Y».
  \item Busca en las opciones la que tenga X y Y, aunque el tema sea otro.
  \item Desconfía de la opción que repite el tema del texto sin el mecanismo.
\end{enumerate}
\end{regla}

\begin{texto}{La paradoja de la elección}{Ensayo de divulgación · texto original escrito para este libro}
\pa Solemos pensar que tener más opciones siempre es mejor. Si una tienda ofrece cinco sabores de mermelada, una que
ofrezca treinta debería atraer más clientes y dejarlos más contentos. Sin embargo, varias observaciones sugieren lo
contrario: cuando las opciones son demasiadas, muchas personas terminan sin elegir nada, o eligen y quedan menos
satisfechas.

\pa La explicación tiene dos partes. La primera es el cansancio: comparar treinta cosas exige un esfuerzo que muchos
prefieren evitar, y la decisión se aplaza. La segunda es el arrepentimiento: con tantas alternativas, es fácil imaginar
que alguna de las descartadas habría sido mejor, y esa sospecha le quita gusto a lo elegido.

\pa Esto no significa que la libertad de elegir sea mala, sino que tiene un punto óptimo. Más allá de él, cada opción nueva
agrega más carga que beneficio. Quienes diseñan menús, formularios o planes de estudio harían bien en recordarlo.
\end{texto}

\subsection{Preguntas comentadas}
\begin{comentada}{¿Cuál situación ilustra mejor la idea del texto?}
\mal{A}{Un restaurante con un menú de cuatro platos tiene pocos clientes porque su comida es mala.}{Comparte el tema
(menú), no el principio: la causa es la calidad, no el exceso de opciones.}
\bien{B}{Un estudiante, ante cincuenta cursos electivos, aplaza la inscripción y, cuando por fin elige, duda de su
decisión.}{Tiene los dos mecanismos: cansancio (aplaza) y arrepentimiento (duda).}
\mal{C}{Una persona elige rápido entre dos camisas.}{Es el caso contrario (pocas opciones), sin conflicto.}
\mal{D}{Una tienda aumenta sus ventas al ampliar su catálogo.}{\emph{Dice lo contrario} del texto.}
\end{comentada}

\taller
\piso{Como en el examen}{opción múltiple}
\begin{ejercicios}
\item Según el texto, ¿qué le recomendaría el autor a quien diseña el menú de una cafetería?
\opcionesl{Ofrecer la mayor cantidad posible de platos.}{Ofrecer una variedad razonable, sin excesos.}{Ofrecer un solo plato.}{No preocuparse por el número de platos.}
\sol{B. Hay un «punto óptimo».}
\item ¿Cuál caso ilustra el segundo mecanismo (el arrepentimiento)?
\opcionesl{Alguien deja de ir a una tienda porque está lejos.}{Alguien compra un celular entre muchos modelos y luego piensa que otro habría sido mejor.}{Alguien no compra nada porque no tiene dinero.}{Alguien regala un libro.}
\sol{B.}
\end{ejercicios}

\begin{retoresuelto}
El principio es «el esfuerzo invertido aumenta el valor percibido». La mesa armada por uno mismo lo cumple: misma mesa,
más esfuerzo, más aprecio. \textbf{Opción A.}
\end{retoresuelto}
\bitacora

\leccion{pu-ap-2}{Analogías}{Juego de las parejas que se relacionan igual.}

\begin{reto}
LIBRO es a BIBLIOTECA como CUADRO es a
\opcionesd{pintor}{museo}{marco}{color}
\end{reto}

\tecnica{La técnica: nombrar la relación}
Una \clave{analogía} compara dos parejas que se relacionan de la misma manera. El método es nombrar la relación de la
primera pareja con una frase precisa («un libro \emph{se guarda y exhibe en} una biblioteca») y aplicar la misma frase a la
segunda. Relaciones frecuentes: parte-todo, objeto-lugar, causa-efecto, herramienta-oficio, sinónimos, antónimos,
grado (tibio-caliente), producto-productor.

\subsubsection{Analogías en los textos}
En lectura, la pregunta suele ser: «la situación del texto es comparable a…». Funciona igual: se describe la relación
central del texto en abstracto y se busca la opción con la misma relación.

\figura{analogia}{En una analogía se repite la relación, no las cosas.}

\begin{regla}
\begin{itemize}
  \item Escribe la relación como una oración: «A \_\_\_ B».
  \item Si varias opciones encajan, haz la frase más precisa.
  \item Respeta el orden: A es a B como C es a D (no D a C).
\end{itemize}
\end{regla}

\begin{texto}{El jardinero y el escultor}{Ensayo breve · texto original escrito para este libro}
\pa Hay dos maneras de pensar la educación. La primera imagina al maestro como un escultor: tiene en la cabeza la forma
final y quita de la piedra todo lo que sobra hasta obtenerla. El estudiante es materia que se moldea.

\pa La segunda imagina al maestro como un jardinero. El jardinero no fabrica la planta: prepara la tierra, riega, poda y
protege de las plagas, pero la planta crece por su cuenta y no siempre como él esperaba. Su trabajo es crear condiciones,
no imponer formas.

\pa Ningún maestro es solo una cosa. Pero la imagen que cada uno elige cambia lo que hace en el salón: el escultor corrige;
el jardinero acompaña.
\end{texto}

\subsection{Preguntas comentadas}
\begin{comentada}{Según el texto, el maestro-escultor es al estudiante como}
\mal{A}{el jardinero es a la planta.}{Es la otra imagen, con la relación contraria (crear condiciones).}
\bien{B}{el alfarero es al barro.}{Quien da forma a una materia según un modelo previo: la misma relación.}
\mal{C}{el lector es al libro.}{No hay acción de dar forma.}
\mal{D}{la semilla es al árbol.}{Relación de origen, no de moldear.}
\end{comentada}

\taller
\piso{Calienta}{relaciones}
\begin{ejercicios}
\item Completa: TIJERA es a SASTRE como BISTURÍ es a \_\_\_. \sol{cirujano (herramienta-oficio).}
\item Completa: SEMILLA es a ÁRBOL como HUEVO es a \_\_\_. \sol{ave (origen-ser adulto).}
\end{ejercicios}
\piso{Como en el examen}{opción múltiple}
\begin{ejercicios}
\item TIBIO es a HIRVIENTE como
\opcionesd{frío es a helado}{nieve es a hielo}{calor es a sol}{agua es a vapor}
\sol{A. Relación de grado: lo moderado frente a lo extremo.}
\item El papel del maestro-jardinero se parece más al de
\opcionesl{un entrenador que diseña rutinas según lo que cada deportista necesita.}{un fabricante que produce piezas idénticas.}{un juez que castiga las faltas.}{un vendedor.}
\sol{A. Crea condiciones y acompaña el desarrollo.}
\end{ejercicios}

\begin{retoresuelto}
Un libro se guarda y exhibe en una biblioteca; un cuadro, en un museo. \textbf{Opción B.} «Pintor» sería la relación
producto-productor, que no es la de la primera pareja.
\end{retoresuelto}
\bitacora

\evaluacion
\begin{texto}{El efecto de la ventana rota}{Divulgación · texto original escrito para este libro}
\pa Una idea muy discutida en el urbanismo sostiene que el descuido invita al descuido. Si en un edificio aparece una
ventana rota y nadie la arregla, pronto aparecerán otras; si en una esquina se acumula basura, más gente la usará como
botadero. La señal que reciben las personas es simple: «aquí a nadie le importa».

\pa La idea tiene defensores y críticos. Los primeros señalan que arreglar rápido los daños pequeños puede prevenir
problemas mayores. Los segundos advierten que se ha usado para justificar medidas excesivas contra faltas menores, y que
el descuido de un barrio suele tener causas más profundas, como la pobreza o el abandono estatal.

\pa Lo que pocos discuten es la parte más modesta del argumento: el entorno comunica normas. Un espacio cuidado sugiere
que hay alguien pendiente; uno abandonado, que no lo hay.
\end{texto}
\begin{ejercicios}
\item ¿Cuál situación ilustra mejor la idea del \S\,1?
\opcionesl{En un parque recién pintado, los visitantes recogen su basura.}{En un salón con las paredes rayadas y sin limpiar, los estudiantes rayan más.}{Un barrio pobre recibe ayuda del Estado.}{Un edificio nuevo tiene ventanas grandes.}
\sol{B. Descuido que invita a más descuido. (2.1)}
\item La situación A de la pregunta anterior ilustra
\opcionesl{la idea contraria a la del texto.}{la «parte más modesta» del argumento: un entorno cuidado comunica que hay normas.}{la crítica a la idea.}{un caso de pobreza.}
\sol{B. (2.1)}
\item VENTANA ROTA es a DESCUIDO como
\opcionesd{fiebre es a enfermedad}{casa es a ventana}{basura es a calle}{arreglo es a daño}
\sol{A. Señal visible de un estado. (2.2)}
\item Según los críticos, la idea es problemática porque
\opcionesl{las ventanas no se rompen solas.}{se ha usado para justificar medidas excesivas y oculta causas más profundas.}{nadie arregla las ventanas.}{el urbanismo no estudia barrios.}
\sol{B. (2.1)}
\end{ejercicios}
